\Watermark    
\begin{center}
    \begin{minipage}[t]{0.2\linewidth}
        % Content for the first part (left logo)
        \begin{center}
            \includegraphics[scale=.5]{images/schools/moey.png}
        \end{center}
    \end{minipage}%
    \hfill
    \begin{minipage}[t]{0.6\linewidth}
        % Content for the second part (middle text)
        \vspace{-3.2cm} % Adjust the space as needed
        \begin{center}
            \fontsize{18pt}{15pt} \textbf{\khmerfontmoul ក្រសួងអប់រំ យុវជន និងកីឡា}\\
            \vspace{0.5cm}
            \fontsize{17pt}{15pt} \textbf{\khmerfontmoul វិទ្យាស្ថានបច្ចេកវិទ្យាកម្ពុជា}
        \end{center}
    \end{minipage}%
    \hfill
    \begin{minipage}[t]{0.2\linewidth}
        % Content for the third part (right logo)
        \begin{center}
            \vspace{-3.4cm}
            \includegraphics[scale=.31]{images/schools/itc.png}
        \end{center}
    \end{minipage}
\end{center}

\begin{center}
    \bfseries 
    \vspace{-.5cm}
    \fontsize{14pt}{15pt} {\khmerfontmoul ដេប៉ាតឺម៉ង់គណិតវិទ្យាអនុវត្តន៍ និងស្ថិតិ} \\[1.5cm]
    %\fontsize{14pt}{15pt} {\khmerfontmoul គម្រោងសញ្ញាប័ត្រវិស្វករ}\\[0.5cm]
    \begin{flushleft}
        \hspace{3.5cm}
        \fontsize{14pt}{15pt} {\khmerfontmoul គម្រោងសញ្ញាប័ត្រវិស្វករ}\\[0.5cm]
        % Student / defense-date rows: label column auto-sizes to its widest
        % entry, so the colons share one fixed column (reference-thesis style).
        {\fontsize{13pt}{18pt}\selectfont
        \begin{tabular}{@{}>{\khmerfontmoul}l@{\,:\hspace{0.55cm}}>{\khmerfontmoul\raggedright\arraybackslash}p{8.2cm}@{}}
            របស់និស្សិត & \emph{\href{mailto:e20210684@itc.edu.kh}{លុន ច័ន្ទប៉ូលី}} \\[0.35cm]
            កាលបរិច្ឆេទការពារនិក្ខេបបទ & ថ្ងៃទី ០៩ ខែកក្កដា ឆ្នាំ ២០២៦ \\
        \end{tabular}}\\[0.8cm]
        \hspace{3.5cm}
        \fontsize{14pt}{15pt} {\khmerfontmoul អនុញ្ញាតឲ្យការពារគម្រោង} \\[1cm]
        {\fontsize{14pt}{18pt}\selectfont
        \begin{tabular}{@{}>{\khmerfontmoul}l@{\,:\hspace{0.55cm}}l@{}}
            នាយកវិទ្យាស្ថាន & \underline{\hspace{6.5cm}} \\
        \end{tabular}}\\[1cm]
        \hspace{3.5cm}
        \fontsize{13pt}{15pt} {\khmerfontmoul រាជធានីភ្នំពេញ, ថ្ងៃទី ០៩ ខែ កក្កដា ឆ្នាំ ២០២៦} \\[0.7cm]
        % Title row: colon hugs the label; the long Khmer title wraps with a
        % hanging indent under its own first line (reference-thesis style).
        {\fontsize{13pt}{18pt}\selectfont
        \begin{tabular}{@{}>{\khmerfontmoul}l@{\,:\hspace{0.4cm}}>{\khmerfontmoul\raggedright\arraybackslash}p{12cm}@{}}
            ប្រធានបទ & ការវាយតម្លៃ និងទទួលធានា (Underwriting) បែបបត់បែន ផ្នែកធានារ៉ាប់រងសុខភាព តាមរយៈ Contextual Bandits៖ វិធីសាស្ត្រនៃ Reinforcement Learning សម្រាប់ប្រទេសកម្ពុជា \\
        \end{tabular}}\\[0.5cm]
        % Establishment / staff rows: shared fixed colon column as above.
        {\fontsize{13pt}{18pt}\selectfont
        \begin{tabular}{@{}>{\khmerfontmoul}l@{\,:\hspace{0.55cm}}>{\khmerfontmoul\raggedright\arraybackslash}p{8cm}@{}}
            សហគ្រាស & ក្រុមហ៊ុន Decent Actuarial Consultants Co., Ltd.\ (DAC) \\[0.35cm]
            ប្រធានដេប៉ាតឺម៉ង់ & បណ្ឌិត លិន មង្គលសិរី \\[0.35cm]
            សាស្ត្រាចារ្យណែនាំ & បណ្ឌិត ហាស់ សុធា \\[0.35cm]
            អ្នកទទួលខុសត្រូវក្នុងសហគ្រាស & លោក អន រ៉ាដែត \\
        \end{tabular}}\\[1cm]
    \end{flushleft}

    \fontsize{13pt}{15pt}{\khmerfontmoul រាជធានីភ្នំពេញ, ២០២៦}

\end{center}